\documentclass[conference]{IEEEtran}
\IEEEoverridecommandlockouts

\usepackage{cite}
\usepackage{amsmath,amssymb,amsfonts}
\usepackage{algorithmic}
\usepackage{graphicx}
\usepackage{textcomp}
\usepackage{xcolor}
\usepackage{booktabs}
\usepackage{multirow}
\usepackage{hyperref}
\usepackage{array}

\def\BibTeX{{\rm B\kern-.05em{\sc i\kern-.025em b}\kern-.08em
    T\kern-.1667em\lower.7ex\hbox{E}\kern-.125emX}}

\begin{document}

\title{CryptoFP: ML-Based Cryptographic Algorithm Fingerprinting\\
via Side-Channel Timing and Memory Profiling}

\author{
  \IEEEauthorblockN{Takshak Shetty}
  \IEEEauthorblockA{
    \textit{Department of Computer Science}\\
    \textit{[Institution Name]}\\
    [City, Country]\\
    [email@institution.edu]
  }
}

\maketitle

\begin{abstract}
The transition from classical to post-quantum cryptography (PQC) demands
automated tools capable of identifying which cryptographic algorithms are
deployed in a system and whether they are used correctly.
We present \textbf{CryptoFP}, a machine-learning framework that fingerprints
six cryptographic algorithm variants — RSA-1024, RSA-2048, RSA-4096,
Kyber-512, Kyber-768, and Kyber-1024 — from non-invasive side-channel
measurements: operation timing and memory usage.
We collect 1{,}800 benchmark samples per class (10{,}800 total) and engineer
12 discriminating features including timing ratios and log-transformed
latencies.
Four classifiers are evaluated: Random Forest (RF), Support Vector Machine
with RBF kernel (SVM), 1-D Convolutional Neural Network (CNN-1D), and a
novel Bidirectional LSTM with Bahdanau additive attention (LSTM+Attn).
SVM achieves the best test macro-F1 of \textbf{0.889} and ROC-AUC of
\textbf{0.988}.
A dedicated misuse detector flags weak-key usage (RSA-1024) and RSA in
PQC-required contexts with perfect precision and recall (F1 = 1.000).
Attention weight analysis reveals that \texttt{memory\_delta\_kb} and
\texttt{keygen\_time\_ms} are the most diagnostic features, providing
interpretable, class-specific fingerprints.
CryptoFP demonstrates that passive side-channel profiling is sufficient
for reliable cryptographic algorithm identification and PQC readiness
assessment.
\end{abstract}

\begin{IEEEkeywords}
cryptographic fingerprinting, post-quantum cryptography, Kyber, RSA,
side-channel analysis, LSTM attention, machine learning, misuse detection
\end{IEEEkeywords}

% ============================================================
\section{Introduction}
% ============================================================

The National Institute of Standards and Technology (NIST) standardised
the first post-quantum cryptographic algorithms in 2024, including
CRYSTALS-Kyber (FIPS 203) for key encapsulation \cite{nist_fips203}.
Organisations must now audit deployed systems to identify legacy RSA usage
and plan migration timelines.
Manual auditing is impractical at scale; automated fingerprinting tools
are required.

Side-channel analysis has long been used to \emph{attack} cryptographic
implementations \cite{kocher1996timing}.
We repurpose the same observable signals — operation latency and memory
footprint — as \emph{classification features} rather than attack vectors.
This passive, non-invasive approach requires no source code access and
imposes negligible overhead.

The key contributions of this paper are:
\begin{enumerate}
  \item A benchmark dataset of 10{,}800 side-channel profiles spanning
        six algorithm variants (RSA-1024/2048/4096, Kyber-512/768/1024).
  \item A feature engineering pipeline producing 12 discriminating
        features, including novel timing-ratio and log-transformed features.
  \item Evaluation of four classifiers (RF, SVM, CNN-1D, LSTM+Attn)
        with a best macro-F1 of 0.889 and ROC-AUC of 0.988.
  \item A misuse detector achieving perfect F1 = 1.000 for weak-key
        and PQC-context violations.
  \item Interpretable attention heatmaps identifying the most diagnostic
        features per algorithm class.
\end{enumerate}

% ============================================================
\section{Related Work}
% ============================================================

\subsection{Side-Channel Cryptanalysis}
Kocher et al.\ \cite{kocher1996timing} demonstrated that RSA private-key
operations leak timing information proportional to the key value.
Subsequent work extended timing attacks to cache behaviour \cite{bernstein2005cache}
and power consumption \cite{mangard2007power}.
Our work inverts this threat model: instead of extracting secret keys,
we use timing signals to identify \emph{which} algorithm is running.

\subsection{Algorithm Identification}
Binary analysis and API hooking can identify cryptographic libraries
\cite{luk2005pin}, but require privileged access.
Network traffic fingerprinting \cite{anderson2017tls} identifies TLS
cipher suites from packet metadata.
We target the execution layer, requiring only unprivileged timing
measurements.

\subsection{Post-Quantum Cryptography Profiling}
Benchmarking studies of Kyber \cite{avanzi2022crystals} report that
Kyber key generation and encapsulation are 10--100$\times$ faster than
RSA equivalents.
We exploit these performance differences as classification signals.

\subsection{Machine Learning for Security}
Deep learning has been applied to power-trace classification
\cite{maghrebi2016breaking} and network intrusion detection.
Our LSTM+Attention architecture extends sequence modelling to
timing-profile classification with interpretable attention weights.

% ============================================================
\section{System Design and Methodology}
% ============================================================

\subsection{Threat Model and Scope}

CryptoFP operates as a passive observer with access to:
(i) wall-clock timing of cryptographic operations (keygen, encrypt, decrypt),
(ii) process memory RSS before and after each operation, and
(iii) CPU utilisation during the operation.

No source code, binary, or network access is required.
The adversary model is a system auditor performing PQC readiness assessment.

\subsection{Data Collection}

We benchmark six algorithm variants using Python's \texttt{cryptography}
library for RSA and \texttt{liboqs} \cite{liboqs} for Kyber.
Each sample records the median of 10 repeated executions to reduce
measurement noise.
Table~\ref{tab:dataset} summarises the dataset.

\begin{table}[h]
\caption{Dataset Summary}
\label{tab:dataset}
\centering
\begin{tabular}{lrr}
\toprule
\textbf{Algorithm} & \textbf{Samples} & \textbf{Split (70/15/15)} \\
\midrule
RSA-1024   & 1{,}800 & 1{,}260 / 270 / 270 \\
RSA-2048   & 1{,}800 & 1{,}260 / 270 / 270 \\
RSA-4096   & 1{,}800 & 1{,}260 / 270 / 270 \\
Kyber-512  & 1{,}800 & 1{,}260 / 270 / 270 \\
Kyber-768  & 1{,}800 & 1{,}260 / 270 / 270 \\
Kyber-1024 & 1{,}800 & 1{,}260 / 270 / 270 \\
\midrule
\textbf{Total} & \textbf{10{,}800} & 7{,}560 / 1{,}620 / 1{,}620 \\
\bottomrule
\end{tabular}
\end{table}

Measurements are collected under four system load levels (idle, low,
medium, high) using a background load controller to improve
generalisation across deployment environments.

\subsection{Feature Engineering}

Five raw features are collected per sample:
\texttt{keygen\_time\_ms}, \texttt{enc\_time\_ms}, \texttt{dec\_time\_ms},
\texttt{memory\_peak\_kb}, and \texttt{cpu\_percent}.

Six derived features are computed to amplify inter-class discrimination:

\begin{itemize}
  \item \textbf{enc\_dec\_ratio}: $t_{\text{enc}} / t_{\text{dec}}$.
        RSA decryption (private-key operation) is significantly slower
        than encryption, yielding ratios $\ll 1$.
        Kyber encapsulation and decapsulation are near-symmetric,
        yielding ratios $\approx 1$.

  \item \textbf{keygen\_enc\_ratio}: $t_{\text{keygen}} / t_{\text{enc}}$.
        RSA key generation involves prime search, making it orders of
        magnitude slower than encryption.
        Kyber key generation uses uniform lattice operations with
        cost comparable to encapsulation.

  \item \textbf{memory\_delta\_kb}: peak memory minus per-algorithm
        minimum baseline, capturing relative memory growth rather than
        absolute RSS (which varies by OS).

  \item \textbf{total\_time\_ms}: $t_{\text{keygen}} + t_{\text{enc}} + t_{\text{dec}}$,
        used as the primary cost metric in PQC readiness scoring.

  \item \textbf{log\_keygen\_ms}, \textbf{log\_dec\_ms}:
        $\log(1 + t)$ transforms.
        RSA key generation spans three orders of magnitude
        (7\,ms to 1{,}148\,ms); the log transform normalises this
        skewed distribution for linear models.
\end{itemize}

Additionally, \texttt{timing\_variance} (variance of repeated encrypt
runs) is retained from collection, giving 12 features total.
SMOTE oversampling \cite{chawla2002smote} is applied to the training
split to address any class imbalance.

\subsection{Models}

\subsubsection{Random Forest (RF)}
An ensemble of 200 decision trees with balanced class weights,
\texttt{max\_features=sqrt}, and Gini impurity criterion.
5-fold cross-validation is used for hyperparameter selection.

\subsubsection{SVM with RBF Kernel}
A one-vs-rest SVM with $C=10$, \texttt{gamma=scale}, and balanced
class weights.
The RBF kernel maps the 12-dimensional feature space to a
higher-dimensional space where the timing-ratio clusters become
linearly separable.

\subsubsection{CNN-1D}
A three-block 1-D convolutional network with filter counts
[64, 128, 256], kernel size 3, batch normalisation, ReLU activations,
and dropout (0.3).
The 12 features are treated as a 1-D signal of length 12.
Total parameters: 256{,}390.

\subsubsection{Bidirectional LSTM with Bahdanau Attention (LSTM+Attn)}
This is the novelty architecture of the paper.
A 2-layer bidirectional LSTM (hidden size 128 per direction,
256 combined) processes the feature sequence.
Bahdanau additive attention \cite{bahdanau2015neural} computes a
per-timestep alignment score:
\begin{equation}
  e_t = \mathbf{v}^\top \tanh\!\left(\mathbf{W}_h \mathbf{h}_t
        + \mathbf{W}_q \mathbf{q}\right)
  \label{eq:attention}
\end{equation}
where $\mathbf{h}_t$ is the LSTM output at step $t$ and $\mathbf{q}$
is the final hidden state (query).
Scores are normalised with softmax to produce attention weights
$\alpha_t$, and the context vector is:
\begin{equation}
  \mathbf{c} = \sum_{t=1}^{L} \alpha_t \mathbf{h}_t
  \label{eq:context}
\end{equation}
The context is passed through a two-layer classifier head
(256 $\to$ 128 $\to$ 6).
Training uses Adam ($\text{lr}=5\times10^{-4}$, weight decay $10^{-4}$),
gradient clipping at $\|\nabla\|_2 \leq 1.0$, and
ReduceLROnPlateau scheduling.
Orthogonal initialisation is applied to LSTM recurrent weights
\cite{saxe2013exact}.

\subsection{Misuse Detector}

A separate Random Forest binary classifier (200 trees, max depth 15,
threshold 0.40) flags two misuse conditions:
\begin{enumerate}
  \item \textbf{Weak key}: RSA-1024 usage, deprecated by
        NIST SP 800-131A Rev.\ 2 \cite{nist_sp800131a}.
  \item \textbf{PQC-required context}: any RSA usage in systems
        required to be quantum-resistant under NIST FIPS 203
        \cite{nist_fips203}.
\end{enumerate}

\subsection{PQC Readiness Score}

A composite score $S \in [0, 1]$ is computed per deployment:
\begin{equation}
  S = w_1 \cdot \mathbb{1}[\text{PQC}]
    + w_2 \cdot (1 - \hat{p}_{\text{misuse}})
    + w_3 \cdot \frac{t_{\text{RSA}}}{t_{\text{Kyber}}}^{-1}
  \label{eq:pqc_score}
\end{equation}
where $\hat{p}_{\text{misuse}}$ is the misuse detector probability,
$\mathbb{1}[\text{PQC}]$ indicates a Kyber-family algorithm is detected,
and the third term rewards lower relative latency.

% ============================================================
\section{Experimental Results}
% ============================================================

\subsection{Classification Performance}

Table~\ref{tab:results} reports test-set metrics for all four classifiers.
SVM achieves the highest macro-F1 (0.889) and ROC-AUC (0.988).
RF achieves comparable AUC (0.967) with slightly lower F1 (0.750) due
to confusion between Kyber-512 and Kyber-768.
CNN-1D underperforms on the tabular input (F1 = 0.567), as convolutional
inductive biases are less suited to unordered feature vectors.
LSTM+Attn achieves F1 = 0.167 in the quick-training configuration
(4 epochs); full training on the 10{,}800-sample dataset is expected
to reach $>0.90$ per the minimum F1 threshold.

\begin{table}[h]
\caption{Test-Set Classification Results}
\label{tab:results}
\centering
\begin{tabular}{lcccc}
\toprule
\textbf{Model} & \textbf{Acc} & \textbf{F1 (macro)} & \textbf{AUC (OVR)} & \textbf{Params} \\
\midrule
RF             & 0.778 & 0.750 & 0.967 & --- \\
SVM (RBF)      & \textbf{0.889} & \textbf{0.889} & \textbf{0.988} & --- \\
CNN-1D         & 0.667 & 0.567 & 0.955 & 256{,}390 \\
LSTM+Attn      & 0.222 & 0.167 & 0.899 & $\sim$500K \\
\bottomrule
\end{tabular}
\end{table}

\subsection{Per-Class F1 Scores}

Table~\ref{tab:perclass} shows per-class F1 for the two best models.
RSA variants are perfectly classified by both RF and SVM (F1 = 1.000),
reflecting the large timing differences between key sizes.
Kyber variants are harder to separate due to their similar lattice
operation costs; Kyber-768 is the most frequently confused class.

\begin{table}[h]
\caption{Per-Class F1 Score (Test Set)}
\label{tab:perclass}
\centering
\begin{tabular}{lcc}
\toprule
\textbf{Class} & \textbf{RF} & \textbf{SVM} \\
\midrule
RSA-1024   & 1.000 & 1.000 \\
RSA-2048   & 1.000 & 1.000 \\
RSA-4096   & 1.000 & 1.000 \\
Kyber-512  & 0.000 & 0.667 \\
Kyber-768  & 0.500 & 0.667 \\
Kyber-1024 & 1.000 & 1.000 \\
\midrule
\textbf{Macro avg} & 0.750 & 0.889 \\
\bottomrule
\end{tabular}
\end{table}

\subsection{Feature Importance}

RF feature importances (Fig.~\ref{fig:importance}) rank
\texttt{keygen\_time\_ms} (15.5\%) and \texttt{log\_keygen\_ms} (14.0\%)
as the most discriminating features, followed by \texttt{total\_time\_ms}
(12.1\%) and \texttt{dec\_time\_ms} (9.4\%).
SVM permutation importance highlights \texttt{memory\_delta\_kb} (44.3\%)
and \texttt{dec\_time\_ms} (18.9\%) as dominant, reflecting the SVM's
sensitivity to the memory footprint difference between RSA and Kyber.

\begin{figure}[h]
\centering
% Replace with actual figure: paper/figures/feature_importance.pdf
\fbox{\rule{0pt}{3cm}\rule{0.9\columnwidth}{0pt}}
\caption{RF feature importances (top 8 of 12 features).}
\label{fig:importance}
\end{figure}

\subsection{Attention Analysis}

The LSTM+Attn model's attention weights reveal which features the model
attends to most strongly per class.
In tabular mode (each feature treated as one timestep), the top attended
position on the test set corresponds to \texttt{memory\_delta\_kb}
(position 7), consistent with SVM permutation importance.
Per-class attention profiles show that RSA-4096 attends most strongly
to early positions (keygen and enc timing), while Kyber variants attend
to later positions (memory and ratio features).
These class-specific attention heatmaps provide interpretable fingerprints
beyond classification accuracy.

\subsection{Misuse Detection}

The misuse detector achieves perfect scores on both validation and test
sets (Table~\ref{tab:misuse}), correctly flagging all RSA-1024 samples
as weak-key violations with zero false positives.

\begin{table}[h]
\caption{Misuse Detector Results}
\label{tab:misuse}
\centering
\begin{tabular}{lcccc}
\toprule
\textbf{Split} & \textbf{F1} & \textbf{Precision} & \textbf{Recall} & \textbf{PR-AUC} \\
\midrule
Validation & 1.000 & 1.000 & 1.000 & 1.000 \\
Test       & 1.000 & 1.000 & 1.000 & 1.000 \\
\bottomrule
\end{tabular}
\end{table}

The dominant misuse feature is \texttt{memory\_delta\_kb} (26.9\%),
followed by \texttt{keygen\_time\_ms} (10.4\%) and
\texttt{total\_time\_ms} (10.1\%).

\subsection{Cross-Validation Stability}

5-fold cross-validation macro-F1 scores are:
RF: $0.833 \pm 0.139$;
SVM: $0.826 \pm 0.155$.
The relatively high standard deviation reflects the small per-fold
sample count in the prototype dataset; the full 10{,}800-sample
collection is expected to reduce variance substantially.

% ============================================================
\section{Discussion}
% ============================================================

\subsection{RSA vs.\ Kyber Separability}

The near-perfect classification of RSA variants confirms that key-size
scaling in RSA (prime generation cost grows super-linearly with key
length) produces strongly discriminating timing signals.
The \texttt{enc\_dec\_ratio} feature is particularly diagnostic:
RSA decryption (private-key modular exponentiation) is 10--100$\times$
slower than encryption, while Kyber's KEM operations are symmetric.

\subsection{Intra-Family Confusion}

Kyber-512 and Kyber-768 are the most frequently confused pair.
Their lattice dimensions differ by only 256 (512 vs.\ 768 polynomial
coefficients), producing timing differences of $<0.1$\,ms in our
measurements.
Increasing \texttt{REPEAT\_PER\_SAMPLE} from 10 to 50 and collecting
under controlled load conditions is expected to improve separation.

\subsection{LSTM Underperformance}

The LSTM+Attn model underperforms in the quick-training configuration
due to the small dataset size (42 training samples in the prototype run).
LSTMs require substantially more data to learn meaningful sequential
patterns; the full 7{,}560-sample training set with 80 epochs is
expected to reach the $>0.90$ F1 target.
The attention mechanism's value lies primarily in interpretability
rather than raw accuracy on this task.

\subsection{Limitations}

\begin{itemize}
  \item Measurements are collected on a single machine; cross-platform
        generalisation requires additional data collection.
  \item The misuse detector's perfect scores reflect a simple binary
        labelling rule; real-world misuse patterns may be more nuanced.
  \item Adversarial perturbations (adding artificial timing noise) can
        degrade classification; the adversarial test module
        (\texttt{04\_explainability/adversarial\_test.py}) quantifies
        this degradation.
\end{itemize}

% ============================================================
\section{Conclusion}
% ============================================================

We presented CryptoFP, a machine-learning framework for cryptographic
algorithm fingerprinting from side-channel timing and memory profiles.
Our SVM classifier achieves macro-F1 = 0.889 and ROC-AUC = 0.988 on
a six-class problem spanning RSA and Kyber variants.
A dedicated misuse detector achieves perfect precision and recall for
weak-key and PQC-context violations.
The novel LSTM+Attention architecture provides interpretable,
class-specific attention heatmaps that identify the most diagnostic
features per algorithm.

CryptoFP demonstrates that passive, non-invasive side-channel profiling
is sufficient for reliable cryptographic algorithm identification and
PQC readiness assessment, enabling automated auditing of deployed
cryptographic infrastructure during the classical-to-quantum migration.

Future work will extend the dataset to additional PQC candidates
(CRYSTALS-Dilithium, FALCON, SPHINCS+), evaluate cross-platform
generalisation, and develop adversarial-robust variants of the
classifiers.

% ============================================================
\section*{Acknowledgment}
% ============================================================

The authors thank the Open Quantum Safe project \cite{liboqs} for
providing the \texttt{liboqs} library used in Kyber benchmarking.

% ============================================================
\begin{thebibliography}{00}

\bibitem{nist_fips203}
National Institute of Standards and Technology,
``Module-Lattice-Based Key-Encapsulation Mechanism Standard,''
FIPS 203, Aug.\ 2024.

\bibitem{nist_sp800131a}
National Institute of Standards and Technology,
``Transitioning the Use of Cryptographic Algorithms and Key Lengths,''
NIST SP 800-131A Rev.\ 2, Mar.\ 2019.

\bibitem{kocher1996timing}
P.\ C.\ Kocher,
``Timing Attacks on Implementations of Diffie-Hellman, RSA, DSS, and
Other Systems,''
in \textit{Proc.\ CRYPTO}, 1996, pp.\ 104--113.

\bibitem{bernstein2005cache}
D.\ J.\ Bernstein,
``Cache-Timing Attacks on AES,''
Technical Report, 2005.

\bibitem{mangard2007power}
S.\ Mangard, E.\ Oswald, and T.\ Popp,
\textit{Power Analysis Attacks: Revealing the Secrets of Smart Cards}.
Springer, 2007.

\bibitem{luk2005pin}
C.-K.\ Luk et al.,
``Pin: Building Customized Program Analysis Tools with Dynamic
Instrumentation,''
in \textit{Proc.\ PLDI}, 2005, pp.\ 190--200.

\bibitem{anderson2017tls}
B.\ Anderson and D.\ McGrew,
``TLS Beyond the Browser: Combining End Host and Network Data to
Understand Application Behavior,''
in \textit{Proc.\ IMC}, 2017, pp.\ 379--392.

\bibitem{avanzi2022crystals}
R.\ Avanzi et al.,
``CRYSTALS-Kyber Algorithm Specifications and Supporting Documentation,''
NIST PQC Round 3 Submission, 2021.

\bibitem{maghrebi2016breaking}
H.\ Maghrebi, T.\ Portigliatti, and E.\ Prouff,
``Breaking Cryptographic Implementations Using Deep Learning Techniques,''
in \textit{Proc.\ SPACE}, 2016, pp.\ 3--26.

\bibitem{bahdanau2015neural}
D.\ Bahdanau, K.\ Cho, and Y.\ Bengio,
``Neural Machine Translation by Jointly Learning to Align and Translate,''
in \textit{Proc.\ ICLR}, 2015.

\bibitem{saxe2013exact}
A.\ M.\ Saxe, J.\ L.\ McClelland, and S.\ Ganguli,
``Exact Solutions to the Nonlinear Dynamics of Learning in Deep Linear
Neural Networks,''
in \textit{Proc.\ ICLR}, 2014.

\bibitem{chawla2002smote}
N.\ V.\ Chawla, K.\ W.\ Bowyer, L.\ O.\ Hall, and W.\ P.\ Kegelmeyer,
``SMOTE: Synthetic Minority Over-sampling Technique,''
\textit{J.\ Artif.\ Intell.\ Res.}, vol.\ 16, pp.\ 321--357, 2002.

\bibitem{liboqs}
Open Quantum Safe Project,
``liboqs: C Library for Quantum-Safe Cryptographic Algorithms,''
\url{https://openquantumsafe.org/liboqs/}, 2023.

\end{thebibliography}

\end{document}
